% PAGES: 199-200
% Continues directly after sec06_c3_3.tex (folio 182), which ends mid-sentence
% on the unnumbered display "a_i + b_i x_{i-1} = q_{i-1}, ...; a_i + b_i x_i =
% r_i, ..." -- the sentence resumes here with "and therefore" leading into
% (6.76)-(6.77). No environment was left open.

and therefore
\begin{align}
a_i &= \frac{q_{i-1}x_i - r_ix_{i-1}}{x_i-x_{i-1}}, \qquad \forall\, i = 1:n; \\
b_i &= \frac{r_i-q_{i-1}}{x_i-x_{i-1}}, \qquad \forall\, i = 1:n.
\end{align}

Summarizing, given $v_i$ and $w_i$, $i = 0:n$, we obtain $c_i$ and $d_i$, $i = 1:n$,
from (6.72) and (6.73), and use them to compute $r_i$ and $q_{i-1}$, $i = 1:n$, from
(6.74) and (6.75). In turn, $r_i$ and $q_{i-1}$, $i = 1:n$, are used to obtain $a_i$
and $b_i$, $i = 1:n$, from (6.76) and (6.77).

Since $w_0 = w_n = 0$, see (6.66), we are left with finding the values of $w_i$, for
$i = 1:(n-1)$. To do so, we use the condition (6.67) for the continuity of the first
derivative, i.e., $f_i'(x_i) = f_{i+1}'(x_i)$, for all $i = 1:(n-1)$, which can be
written as
\begin{equation}
b_i + 2c_ix_i + 3d_ix_i^2 \;=\; b_{i+1} + 2c_{i+1}x_i + 3d_{i+1}x_i^2, \qquad \forall\,
i = 1:(n-1);
\end{equation}

From (6.77) and using (6.74) and (6.75), we obtain that
\begin{align*}
b_i &= \frac{v_i-v_{i-1}}{x_i-x_{i-1}} - c_i(x_i+x_{i-1}) -
d_i(x_i^2+x_ix_{i-1}+x_{i-1}^2); \\
b_{i+1} &= \frac{v_{i+1}-v_i}{x_{i+1}-x_i} - c_{i+1}(x_{i+1}+x_i) -
d_{i+1}(x_{i+1}^2+x_{i+1}x_i+x_i^2),
\end{align*}
for all $i = 1:(n-1)$, and therefore
\begin{align*}
b_i + 2c_ix_i + 3d_ix_i^2
&= \frac{v_i-v_{i-1}}{x_i-x_{i-1}} + c_i(x_i-x_{i-1}) + d_i(2x_i^2-x_ix_{i-1}-x_{i-1}^2) \\
&= \frac{v_i-v_{i-1}}{x_i-x_{i-1}} + c_i(x_i-x_{i-1}) + d_i(x_i-x_{i-1})(2x_i+x_{i-1}); \\
b_{i+1}+2c_{i+1}x_i+3d_{i+1}x_i^2
&= \frac{v_{i+1}-v_i}{x_{i+1}-x_i} - c_{i+1}(x_{i+1}-x_i) -
d_{i+1}(x_{i+1}^2+x_{i+1}x_i-2x_i^2) \\
&= \frac{v_{i+1}-v_i}{x_{i+1}-x_i} - c_{i+1}(x_{i+1}-x_i) - d_{i+1}(x_{i+1}-x_i)(x_{i+1}+2x_i),
\end{align*}
for all $i = 1:(n-1)$, which can be written using (6.72) and (6.73) in terms of $w_i$,
$i = 0:n$, as follows:
\begin{align}
b_i+2c_ix_i+3d_ix_i^2
&= \frac{v_i-v_{i-1}}{x_i-x_{i-1}} + \frac{w_{i-1}x_i-w_ix_{i-1}}{2} +
\frac{(w_i-w_{i-1})(2x_i+x_{i-1})}{6} \notag \\
&= \frac{v_i-v_{i-1}}{x_i-x_{i-1}} + w_{i-1}\frac{x_i-x_{i-1}}{6} +
w_i\frac{x_i-x_{i-1}}{3}; \\
b_{i+1}+2c_{i+1}x_i+3d_{i+1}x_i^2
&= \frac{v_{i+1}-v_i}{x_{i+1}-x_i} - \frac{w_ix_{i+1}-w_{i+1}x_i}{2} -
\frac{(w_{i+1}-w_i)(2x_{i+1}+x_i)}{6} \notag \\
&= \frac{v_{i+1}-v_i}{x_{i+1}-x_i} - w_i\frac{x_{i+1}-x_i}{3} -
w_{i+1}\frac{x_{i+1}-x_i}{6},
\end{align}
for all $i = 1:(n-1)$.

From (6.79) and (6.80), we obtain that the constraints (6.78) are equivalent to
\[
\frac{v_i-v_{i-1}}{x_i-x_{i-1}} + w_{i-1}\frac{x_i-x_{i-1}}{6} + w_i\frac{x_i-x_{i-1}}{3}
\;=\;
\frac{v_{i+1}-v_i}{x_{i+1}-x_i} - w_i\frac{x_{i+1}-x_i}{3} - w_{i+1}\frac{x_{i+1}-x_i}{6},
\]
which can be written as
\begin{equation}
w_{i-1}\frac{x_i-x_{i-1}}{6} + w_i\frac{x_{i+1}-x_{i-1}}{3} + w_{i+1}\frac{x_{i+1}-x_i}{6}
\;=\;
\frac{v_{i+1}-v_i}{x_{i+1}-x_i} - \frac{v_i-v_{i-1}}{x_i-x_{i-1}},
\end{equation}
for all $i = 1:(n-1)$.

After multiplication by 6, the equations (6.81) can be expressed as the linear system
\begin{equation}
Mw \;=\; z,
\end{equation}
where $w = (w_i)_{i=1:n-1}$ is the vector of unknowns, $z$ is the $(n-1) \times 1$
vector given by
\[
z(i) \;=\; 6\left(\frac{v_{i+1}-v_i}{x_{i+1}-x_i} - \frac{v_i-v_{i-1}}{x_i-x_{i-1}}\right),
\qquad \forall\, i = 1:(n-1),
\]
and $M$ is the $(n-1) \times (n-1)$ matrix given by
\begin{align}
M(i,i) &= 2(x_{i+1}-x_{i-1}), \qquad \forall\, i = 1:(n-1); \\
M(i,i+1) &= x_{i+1}-x_i, \qquad \forall\, i = 1:(n-2); \\
M(i,i-1) &= x_i - x_{i-1}, \qquad \forall\, i = 2:(n-1).
\end{align}

The matrix $M$ is tridiagonal. Note that $M$ is also symmetric, since, from (6.84) and
(6.85), we find that
\[
M(i,i+1) \;=\; x_{i+1}-x_i \;=\; M(i+1,i), \qquad \forall\, i = 1:(n-2).
\]

Moreover, the matrix $M$ is strictly diagonally dominant, since
\begin{align*}
M(1,1) &= 2(x_2-x_0) \;>\; x_2-x_1 \;=\; M(1,2); \\
M(n-1,n-1) &= 2(x_n-x_{n-2}) \;>\; x_{n-1}-x_{n-2} \;=\; M(n-1,n-2);
\end{align*}
\[
|M(i,i-1)| + |M(i,i+1)| \;=\; x_{i+1}-x_{i-1} \;<\; 2(x_{i+1}-x_{i-1}) \;=\; M(i,i),
\]
for all $i = 2:(n-2)$; cf. (6.83--6.85).

From Theorem 5.7, we conclude that the matrix $M$ is a tridiagonal symmetric positive
definite matrix. Thus, the linear system $Mw = z$ from (6.82) can be solved
efficiently using the LU tridiagonal spd solver from Table 6.6, i.e.,
\[
w \;=\; \text{linear\_solve\_LU\_tridiag\_spd}(M,z).
\]

The pseudocode for the efficient implementation of the natural cubic spline
interpolation can be found in Table 6.8.
% This file ends here (end of PDF page 200 / folio 184), right before Table 6.8,
% which appears on PDF page 201 (out of range for this agent, owned by c4). No
% environment is left open at this point.
